\par\noindent\textbf{(1)} 定义一维快速度χ=artanh(v/c)，四速度Uμ=c(coshχ,sinhχ)。恒定固有加速度意味着dχ/dτ=g0/c，初态静止给χ=g0τ/c，实验室速度v(τ)=c tanh(g0τ/c)，始终严格小于c，不能把v直接写成g0τ。
\par\noindent\textbf{(2)} 利用dt/dτ=coshχ、dx/dτ=c sinhχ积分，得到t=(c/g0)sinh(g0τ/c)，x=(c²/g0)[cosh(g0τ/c)−1]。消去τ有(x+c²/g0)²−c²t²=(c²/g0)²，说明为Rindler双曲轨道，固有加速度与实验室坐标加速度并不相同。
\par\noindent\textbf{(3)} 一维理想变质量火箭动量守恒给dχ=−(u/c)d ln m，故从静止质量m0到快速度χ时 ln(m0/m)=cχ/u。等价 m0/m=exp[(c/u)artanh(v/c)]；u≪c且v≪c时才退化为经典Tsiolkovsky关系v≈u ln(m0/m)。
\par\noindent\textbf{(4)} 短固有时间g0τ/c≪1时x≈g0τ²/2、t≈τ，恢复经典运动学；长时v→c但需要χ→∞，因此理想燃料比指数发散。题中恒定加速度需由喷流主动提供，若把同一运动视作外界推动则不能同时套用内生喷气的质量比公式。
